\chapter{Vergleichende Analyse der Einzelsysteme}\label{ch:vergleichendeAnalyse}

Eine isolierte Betrachtung der etablierten Frühwarnsysteme zeigt, dass jede Technologie über spezifische Stärken verfügt, jedoch auch physikalische oder konzeptionelle Beschränkungen aufweist. Dabei ist zu erkennen, dass alle drei Systeme Unterschiede in Bezug auf die Detektionslatenz, die Fehlalarmresistenz und die Robustheit sowie die Skalierbarkeit und die ökonomische Effizienz aufweisen. Daraus wird ersichtlich, dass keine der Technologien alle Anforderungen optimal erfüllt.

Weltraumgestützte Systeme charakterisieren sich durch ihre globale und großflächige Abdeckung. Satelliten, wie beispielsweise \ac{VIIRS}, verfügen über eine hohe Skalierbarkeit und gestatten die Datenerfassung ohne die Notwendigkeit zusätzlicher physischer Infrastruktur am Boden. Sie liefern  Daten mit gleichbleibender Qualität und fungieren als Grundlage für die Verarbeitung von Informationen aus der ganzen Welt. Gemessen an der überwachten Fläche resultiert dies in einer signifikant hohen ökonomischen Effizienz. Eine Analyse der historischen Entwicklung zeigt, dass in diesem Bereiche größten Schwächen zu verzeichnen waren. Brände wurden von den Satelliten im \acs{LEO} meist erst erkannt, wenn sie bereits eine offene Flamme entwickelt hatten und deshalb eine große thermische Signatur aufwiesen. Kommerziell genutzte Satelliten minimieren diese Latenz mittlerweile durch bordeigenes Edge Computing auf wenige Minuten. Dennoch sind die Systeme durch physikalische Barrieren, wie zum Beispiel dichte Wolkendecken, in ihrer Fehlalarmresistenz und damit in ihrer Zuverlässigkeit eingeschränkt~\cite{DZLR2024, Seton2025,VIIRSLexikon}.

Optisch-terrestrische Systeme, wie IQ FireWatch, bieten demgegenüber eine deutlich höhere Detektionslatenz, da sie ihre Umwelt kontinuierlich in Echtzeit überwachen. Der Einsatz einer multispektralen Sensorik in Kombination mit künstlicher Intelligenz führt zu einer erhöhten Fehlalarmresistenz. Dies ist auf die Fähigkeit zurückzuführen, bekannte Störquellen zu erkennen und zu eliminieren. Für eine lückenlose und landesweite Abdeckung aller Waldflächen, müssten jedoch hunderte Türme mit Strom- und Datenanbindung errichtet werden, was die Kosten und den Wartungsaufwand erheblich steigert. Diese Systeme sind dementsprechend weniger skalierbar und ökonomisch weniger effizient~\cite{Gemmingen2022,Winter2021}.

\acs{IoT}-basierte Sensornetzwerke bieten die geringste Detektionslatenz. Systeme wie \enquote{Silvanet} von Dryad Networks sind in der Lage, Brände bereits in der Schwelbrandphase ab einer Größe von einem Quadratmeter zu erkennen. Zudem ermöglicht das eingesetzte Edge Computing eine lokale Datenverarbeitung nahezu in Echtzeit.~\acs{LPWAN} oder das vor allem genutzte \acs{LoRaWAN} bieten dem System außerdem eine hohe Robustheit im Gelände, da sie auch in Gebieten mit schlechter Netzabdeckung zuverlässig Daten übertragen können. Allerdings sind für eine effektive Abdeckung eines Gebietes hunderte Sensoren erforderlich. Dies führt zu einer erheblichen Steigerung der Kosten und somit zu einer starken Einschränkung der ökonomischen Effizienz bei einer flächendeckenden Anwendung~\cite{Chan2024,Kumar2023, Nodes.sh2024}.

\chapter{Identifikation technologischer Lücken}\label{ch:identifikationTechnologischerLuecken}

Basierend auf dem Vergleich anhand der fünf Kriterien wird deutlich, dass die technologischen Lücken nicht primär in der mangelnden Qualität der Systeme liegen, denn es existieren hochmoderne Satelliten, echtzeitfähige Funknetzwerke und hochentwickelte KI-Algorithmen. Vielmehr besteht die Herausforderung in der Kombination der Technologien, um die jeweiligen Stärken zu nutzen und die Schwächen zu kompensieren~\cite{Chan2024}.

Ein wesentlicher Aspekt stellt hier der Konflikt zwischen der Detektionslatenz und der Skalierbarkeit dar. Während \acs{IoT}-Sensoren eine extrem niedrige Latenz bieten, sind sie aufgrund der erforderlichen Anzahl an Sensorknoten und den damit verbundenen Infrastrukturkosten nicht flächendeckend einsetzbar. Satelliten hingegen bieten eine großflächige Abdeckung und skalieren global sowie effizient, können einen Brand aber erst dann wahrnehmen, wenn er groß genug ist, um das Blätterdach thermisch oder optisch zu durchdringen, was die Echtzeiterkennung oftmals limitiert. Optisch-terrestrische Systeme liegen in der Mitte, können aber durch die erforderliche Turm- und Netzinfrastruktur nicht flächendeckend eingesetzt werden~\cite{Allison2016}.

Ein weiterer Konflikt besteht zwischen der Fehlalarmresistenz und der visuellen Verifizierbarkeit. Ein \acs{IoT}-Sensor ist in der Lage, chemische Veränderungen ohne große Verzögerung zu erkennen, jedoch kann ohne eine visuelle Verifizierung nicht gewährleistet werden, dass es sich tatsächlich um einen Brand handelt~\cite{Kumar2023}. Außerdem fehlt durch das reine Vertrauen auf die Sensordaten das genaue Ausmaß der Flammen, die Ausbreitungsrichtung und die Kenntnis über betroffene Zufahrtswege. Im ungünstigsten Szenario führt dies zu einer Alarmierung auf Basis einer ungesicherten Verdachtslage, ohne dasss den Einsätzkräften eine fundierte taktische Lagebeurteilung zugänglich gemacht wird. Da nicht ersichtlich ist, wie groß der Waldbrand bereits ist, kann es passieren, dass zu viele oder zu wenige Einsatzkräfte alarmiert werden, was mit erhöhten Kosten oder verheerenden Verzögerungen bei den Löscharbeiten einhergeht. Satelliten und optisch-terrestrische Systeme können hingegen visuell verifizieren, sind aber durch Wolken oder topografische Hindernisse in ihrer Reaktionsgeschwindigkeit oftmals eingeschränkt~\cite{Bernhard2010}.

Außerdem besteht eine Lücke in der Robustheit der Systeme und der Kommunikationsinfrastruktur. Obwohl \acs{LoRaWAN}-Sensoren am Boden sehr robust und energieeffizient arbeiten, sind Gateways nötig, um die Daten via 4G, 5G oder Kabel an einen Server weiterzuleiten~\cite{Nodes.sh2024}. In sehr abgelegenen Gebieten fehlt diese notwendige Infrastruktur, wodurch die Alarmierung verzögert oder unmöglich wird. Satelliten und optisch-terrestrische Systeme sind hier weniger abhängig von der lokalen Infrastruktur, können aber durch atmosphärische Bedingungen oder topografische Hindernisse beeinträchtigt werden~\cite{VIIRSLexikon}.

\chapter{Ableitung eines multimodalen Konzepts}\label{ch:ableitungmultimodalenKonzep}

Um die identifizierten technologischen Lücken zu schließen und alle definierten Anforderungen vollumfänglich zu erfüllen, bedarf es eines multimodalen Frühwarnsystems, welches die Stärken der einzelnen Technologien kombiniert und die jeweiligen Schwächen kompensiert. Aktuelle Forschungsprojekte und kommerzielle Entwicklungen zeigen, wie terrestrische \acs{IoT}-Sensorik, autonome Flugsysteme und satellitengestützte Erdbeobachtung in einer hybriden Architektur zusammengeführt werden können.

Ein solches System könnte beispielsweise wie folgt aufgebaut sein:
Um den Konflikt zwischen der Robustheit der Systeme und der Kommunikationsinfrastruktur zu lösen, werden zunehmend terrestrische \acs{LPWAN}-Netzwerke mit Satellitentechnologie gekoppelt. Dryad Networks hat hierfür in Kooperation mit dem französischen Satellitenbetreiber Kinéis eine Technologie entwickelt, die die direkte Übertragung von Daten der \acs{LoRaWAN}-Sensoren an Satelliten realisiert. Diese \enquote{Direct-to-Satellite}-Verbindung sorgt dafür, dass Gateways Alarmmeldungen direkt an eine Satellitenkonstellation im Orbit senden, falls kein terrestrisches Mobilfunknetz zur Verfügung steht. Dies gewährleistet eine lückenlose und weltweite Alarmierungskette, selbst in unzugänglichen Gebirgsregionen oder entlang kritischer linearer Infrastrukturen wie Bahn- und Stromtrassen~\cite{DryadNetworks2024}.

Auch für die Lücke zwischen der Fehlalarmresistenz und der visuellen Verifizierbarkeit lässt sich durch eines multimodales Konzept eine Lösung finden. Die Sensoren am Boden sollen als automatische Auslöser (\enquote{Trigger}) agieren und unbemannte Luftfahrzeuge (\acs{UAV}s/Drohnen) zur visuellen Lagebeurteilung anfordern und autonom zum Detektionsort navigieren. Über ein latenzarmes 5G-Netz kann die Drohne so vor dem Eintreffen der ersten Einsatzkräfte hochauflösende optische und thermische Echtzeit-Videodaten an die Einsatzzentrale übermitteln. Die visuelle Verifikation erlaubt eine Einschätzung der Brandausbreitung und der betroffenen Zufahrtswege. Hierdurch wird eine gezielte Alarmierung der Einsatzkräfte gewährleistet und folglich können Kosten und Verzögerungen bei Löscharbeiten minimiert werden~\cite{Dauernheim2025, Laessig2025}.

Moderne Kleinsatelliten-Konstellationen bieten die Lösung für den Konflikt zwischen Detektionslatenz und Skalierbarkeit für Regionen, in denen eine flächendeckende \acs{IoT}-Sensorik unwirtschaftlich ist. Unternehmen wie OroraTech setzen hierbei auf Edge-Computing im Orbit und sogenannte \enquote{Thermal-Livestreams}~\cite{Seton2025}. 
Damit das multimodale Konzept auch praktisch umsetzbar ist, müssen alle Datenströme in einer zentralen Einsatzzentrale zusammengeführt werden. Im Forschungsprojekt \enquote{5G-Waldwächter} passiert das mittels Unterstützung einer KI über einen Zentralen Krisenmanagement-Server (\acs{ZKMS}). Dies ermöglicht auch Voraussagen zu erwarteten Brandausbreitungen, um Einsatzkräfte frühzeitig auf die zu erwartende Lage vorzubereiten~\cite{Laessig2025}. 



